% preamble-ko.inc.tex

\usepackage{kotex}
\usepackage{tikz}
\usepackage{emoji}
\usepackage[absolute,overlay]{textpos}
\usepackage{algorithm}
\usepackage{algpseudocode}
\usepackage{fontawesome}
\usepackage[normalem]{ulem}
\usepackage{etoolbox}

\newfontfamily\lmfont{Latin Modern Roman}
\AtBeginEnvironment{algorithmic}{\lmfont}
\setemojifont{Noto Color Emoji}
\setmainhangulfont{NanumSquareRound}

\usetheme{metropolis}

\definecolor{mainred}{HTML}{B71C1C}
\definecolor{mainblue}{HTML}{1565C0}
\definecolor{maingreen}{HTML}{2E7D32}
\definecolor{mainorange}{HTML}{EF6C00}
\definecolor{maingray}{HTML}{616161}

% Markdown 링크와 직접 작성한 \href에 같은 밑줄 스타일을 적용한다.
% 링크 색상은 Pandoc이 hypersetup을 만들 때 build.py의 변수로 지정한다.
\AtBeginDocument{%
  \NewCommandCopy{\mkslideoriginalhref}{\href}%
  \RenewDocumentCommand{\href}{o m m}{%
    \IfNoValueTF{#1}%
      {\mkslideoriginalhref{#2}{\uline{#3}}}%
      {\mkslideoriginalhref[#1]{#2}{\uline{#3}}}%
  }%
}

% ===== 사용자 정의 명령 =====

% 세로 구분선의 좌우 간격: \columnbar[left=0.5em,right=1em]
\pgfkeys{
  /mkslide/columnbar/.is family,
  /mkslide/columnbar,
  left/.initial=0.5em,
  right/.initial=1em,
}
\NewDocumentCommand{\columnbar}{O{}}{%
  \begingroup
  \pgfkeys{/mkslide/columnbar/.cd,#1}%
  \hspace{\pgfkeysvalueof{/mkslide/columnbar/left}}%
  \vrule
  \hspace{\pgfkeysvalueof{/mkslide/columnbar/right}}%
  \endgroup
}

% 의미 기반 강조: 기본은 bold, [normal] 옵션으로 보통 굵기 지정
\NewDocumentCommand{\mkslideemphasis}{m m m}{%
  \ifstrequal{#1}{bold}{%
    \textcolor{#2}{\textbf{#3}}%
  }{%
    \ifstrequal{#1}{normal}{%
      \textcolor{#2}{#3}%
    }{%
      \PackageError{mkslide}{Unknown emphasis style '#1'}%
        {Use 'bold' or 'normal'.}%
    }%
  }%
}
\NewDocumentCommand{\critical}{O{bold} m}{%
  \mkslideemphasis{#1}{mainred}{#2}%
}
\NewDocumentCommand{\concept}{O{bold} m}{%
  \mkslideemphasis{#1}{mainblue}{#2}%
}
\NewDocumentCommand{\positive}{O{bold} m}{%
  \mkslideemphasis{#1}{maingreen}{#2}%
}
\NewDocumentCommand{\caution}{O{bold} m}{%
  \mkslideemphasis{#1}{mainorange}{#2}%
}
\NewDocumentCommand{\muted}{O{bold} m}{%
  \mkslideemphasis{#1}{maingray}{#2}%
}

% 글자색은 유지하고 밑줄에만 색상 적용
\makeatletter
\NewDocumentCommand{\coloruline}{O{mainred} m}{%
  \leavevmode\bgroup
  \UL@setULdepth
  \markoverwith{%
    \lower\ULdepth\hbox{\color{#1}\rule{0.5em}{\ULthickness}}%
  }%
  \ULon{#2}%
}
\NewDocumentCommand{\coloruwave}{O{mainred} m}{%
  \leavevmode\bgroup
  \ifdim\ULdepth=\maxdimen
    \ULdepth=3.5\p@
  \else
    \advance\ULdepth by 2\p@
  \fi
  \markoverwith{\lower\ULdepth\hbox{\color{#1}\sixly\char58}}%
  \ULon{#2}%
}
\makeatother

% 자주 쓰는 밑줄 색상 단축 명령
\NewDocumentCommand{\reduline}{m}{\coloruline[mainred]{#1}}
\NewDocumentCommand{\blueuline}{m}{\coloruline[mainblue]{#1}}
\NewDocumentCommand{\greenuline}{m}{\coloruline[maingreen]{#1}}
\NewDocumentCommand{\reduwave}{m}{\coloruwave[mainred]{#1}}
\NewDocumentCommand{\blueuwave}{m}{\coloruwave[mainblue]{#1}}
\NewDocumentCommand{\greenuwave}{m}{\coloruwave[maingreen]{#1}}

% 전체 배경/본문
\setbeamercolor{normal text}{fg=black,bg=white}

% metropolis 포인트 컬러(불릿/섹션 등)
\setbeamercolor{structure}{fg=mainred}

% 헤더(제목 영역) 흰색
\setbeamercolor{frametitle}{bg=white, fg=black}
\setbeamerfont{frametitle}{series=\bfseries}

% Footnote font size
\setbeamerfont{footnote}{size=\tiny}

% ===== underbar (progress bar) 유지/색 지정 =====
\metroset{progressbar=frametitle}
\setbeamercolor{progress bar}{fg=mainred}

% 불릿 모양
\setlength{\leftmargini}{1.5em}    % 1레벨 들여쓰기
\setlength{\leftmarginii}{1.2em}   % 2레벨
\setlength{\leftmarginiii}{1.0em}  % 3레벨

\setbeamertemplate{itemize item}{$\bullet$}
\setbeamertemplate{itemize subitem}{$\circ$}
\setbeamertemplate{itemize subsubitem}{$\cdot$}

% block의 indent 옵션에 사용하는 좌우 여백 환경.
% leftskip/rightskip은 itemize/enumerate가 시작될 때 초기화되므로,
% 바깥쪽 list로 본문 폭을 줄여 내부 목록에도 여백이 유지되게 한다.
\makeatletter
\newenvironment{mkslideindent}[1]{%
  \begin{list}{}{%
    \setlength{\leftmargin}{#1}%
    \setlength{\rightmargin}{#1}%
    \setlength{\listparindent}{\parindent}%
    \setlength{\topsep}{0pt}%
    \setlength{\partopsep}{0pt}%
    \setlength{\parsep}{\parskip}%
    \setlength{\itemsep}{0pt}%
  }%
  \item\relax
  % 보이지 않는 바깥 list가 사용자 목록의 중첩 레벨을 소비하지 않게 한다.
  \advance\@listdepth\m@ne
}{%
  \advance\@listdepth\@ne
  \end{list}%
}
\makeatother

% block 스타일
\setbeamercolor{block title}{bg=orange!70, fg=white}
\setbeamercolor{block body}{bg=orange!10, fg=black}

\setbeamercolor{block title alerted}{bg=mainred, fg=white}
\setbeamercolor{block body alerted}{bg=mainred!15, fg=black}

\setbeamercolor{block title example}{bg=green!50!black, fg=white}
\setbeamercolor{block body example}{bg=green!10, fg=black}

% Page Number
\setbeamertemplate{footline}{
  \hfill
  \insertframenumber/\inserttotalframenumber
  \hspace{0.3cm}
  \vspace{0.2cm}
}

% ===== school mark (bottom-left, fixed) =====
\makeatletter
\addtobeamertemplate{background canvas}{}{%
  \ifbeamer@plainframe\relax
  \else
    \ifnum\insertframenumber=1\relax
    \else
      \begin{textblock*}{9mm}(2mm,\paperheight-8mm)%
        \includegraphics[width=9mm]{@@LOGO_PATH@@}%
      \end{textblock*}%
    \fi
  \fi
}
\makeatother

% 제목 없을 때 제목 행은 숨기되 배경은 유지하는 템플릿
\makeatletter
\setbeamertemplate{block begin}{%
  \@ifempty\insertblocktitle{}{%
    \begin{beamercolorbox}[sep=0.5ex,vmode]{block title}%
      \usebeamerfont{block title}\insertblocktitle%
    \end{beamercolorbox}%
    \nointerlineskip%
  }%
  \begin{beamercolorbox}[sep=1ex,vmode]{block body}%
    \usebeamerfont{block body}%
}
\setbeamertemplate{block end}{%
  \end{beamercolorbox}%
}
\makeatother

\tracinglostchars=0
